\documentclass[11pt,a4paper]{article}
\usepackage[utf8]{inputenc}
\usepackage[T1]{fontenc}
\usepackage[spanish]{babel}
\usepackage{geometry}
\usepackage{amsmath}
\usepackage{amssymb}
\usepackage{enumitem}
\usepackage{booktabs}
\usepackage{fancyhdr}
\usepackage{xcolor}
\usepackage{mdframed}
\usepackage{titlesec}
\usepackage{multicol}

\geometry{top=2cm, bottom=2.5cm, left=2.5cm, right=2.5cm}

\pagestyle{fancy}
\fancyhf{}
\rhead{Tecnolog\'ia Digital VI --- 1er Semestre 2026}
\lhead{Banco de Preguntas Te\'oricas --- Solo Preguntas}
\rfoot{P\'agina \thepage}

\definecolor{tdblue}{RGB}{0,74,135}
\definecolor{tdgreen}{RGB}{0,112,60}
\definecolor{tdred}{RGB}{180,35,24}
\definecolor{tdorange}{RGB}{200,100,0}
\definecolor{lightgray}{RGB}{245,245,245}
\definecolor{lightblue}{RGB}{230,240,255}
\definecolor{lightgreen}{RGB}{230,248,235}
\definecolor{lightyellow}{RGB}{255,252,220}

% Caja de sección
\newmdenv[linecolor=tdblue, linewidth=2pt, roundcorner=5pt,
  backgroundcolor=lightblue, innertopmargin=8pt, innerbottommargin=8pt,
  innerleftmargin=10pt]{sectionbox}

% Caja de advertencia (pregunta clásica)
\newmdenv[linecolor=tdred, linewidth=2pt, roundcorner=5pt,
  backgroundcolor=lightyellow, innertopmargin=8pt, innerbottommargin=8pt,
  innerleftmargin=10pt]{classicbox}

\titleformat{\section}{\large\bfseries\color{tdblue}}{}{0em}{}[\titlerule]
\titleformat{\subsection}{\normalsize\bfseries\color{tdblue}}{}{0em}{}

\setlist[enumerate]{itemsep=8pt}

\begin{document}

%--- Portada ---
\begin{center}
  {\Large\bfseries\color{tdblue} Universidad Torcuato Di Tella}\\[6pt]
  {\large\bfseries Tecnolog\'ia Digital VI}\\[4pt]
  {\large Banco de Preguntas Te\'oricas para el Parcial}\\[4pt]
  {\normalsize 1er Semestre 2026}\\[12pt]
  \rule{\linewidth}{1.5pt}
\end{center}

\begin{mdframed}[linecolor=tdorange, linewidth=1.5pt, roundcorner=4pt,
  backgroundcolor=lightyellow]
\textbf{C\'omo usar este documento:} Este banco cubre todos los temas te\'oricos
evaluables en el parcial. Las preguntas marcadas con
\textcolor{tdred}{\textbf{[CL\'ASICA]}} son las que se repiten hist\'oricamente.
Las respuestas se encuentran en el documento \textit{preguntas\_teoricas.pdf}.
\end{mdframed}

\vspace{0.5cm}
\tableofcontents
\newpage

% =============================================================
\section{Preprocesamiento y Validaci\'on}
% =============================================================

\subsection{Valores Faltantes}

\begin{enumerate}[label=\textbf{P\arabic*.}]

\item Explique los tres mecanismos de datos faltantes (MCAR, MAR, MNAR) y d\'e un ejemplo
  de cada uno.

\item ¿Qu\'e estrategias existen para imputar valores faltantes? ¿Cu\'ando usar cada una?

\end{enumerate}

\subsection{Data Leakage}

\begin{enumerate}[label=\textbf{P\arabic*.}, resume]

\item \leavevmode

\begin{classicbox}
\textcolor{tdred}{\textbf{[CL\'ASICA]}} ¿Qu\'e es el \textit{data leakage}? D\'e un
ejemplo concreto y explique c\'omo prevenirlo.
\end{classicbox}

\item Un equipo aplica un MinMaxScaler ajustado sobre todo el dataset y luego divide en
  train/test. ¿Hay data leakage? Justifique.

\end{enumerate}

\subsection{Escalado de Features}

\begin{enumerate}[label=\textbf{P\arabic*.}, resume]

\item ¿Cu\'ando es obligatorio escalar features y cu\'ando no importa? D\'e ejemplos de
  modelos en cada categor\'ia.

\item Compare StandardScaler, MinMaxScaler y RobustScaler. ¿Cu\'ando preferir\'ia cada uno?

\end{enumerate}

\subsection{Desbalance de Clases}

\begin{enumerate}[label=\textbf{P\arabic*.}, resume]

\item \leavevmode

\begin{classicbox}
\textcolor{tdred}{\textbf{[CL\'ASICA]}} Un dataset de detecci\'on de fraude tiene 98\%
de transacciones normales y 2\% fraudulentas. Un modelo que predice siempre ``normal''
obtiene 98\% de accuracy. ¿Qu\'e problema hay y qu\'e har\'ia?
\end{classicbox}

\item ¿Qu\'e es SMOTE y c\'omo genera muestras sint\'eticas?

\end{enumerate}

\subsection{Validaci\'on Cruzada}

\begin{enumerate}[label=\textbf{P\arabic*.}, resume]

\item \leavevmode

\begin{classicbox}
\textcolor{tdred}{\textbf{[CL\'ASICA]}} Explique K-Fold Cross-Validation. ¿Por qu\'e
es preferible al simple holdout?
\end{classicbox}

\item ¿Qu\'e es Leave-One-Out CV (LOOCV)? ¿Cu\'al es su ventaja y su desventaja principal?

\end{enumerate}

% =============================================================
\section{M\'etricas de Evaluaci\'on}
% =============================================================

\subsection{Regresi\'on}

\begin{enumerate}[label=\textbf{P\arabic*.}, resume]

\item Compare MSE, RMSE y MAE. ¿Cu\'ando elegir\'ia cada uno?

\item ¿Puede el $R^2$ ser negativo? Explique qu\'e significar\'ia eso.

\item ¿Cu\'ando usar RMSLE en lugar de RMSE?

\end{enumerate}

\subsection{Clasificaci\'on}

\begin{enumerate}[label=\textbf{P\arabic*.}, resume]

\item \leavevmode

\begin{classicbox}
\textcolor{tdred}{\textbf{[CL\'ASICA]}} Defina Precision y Recall. ¿Cu\'ando es m\'as
grave tener Recall bajo? ¿Cu\'ando es m\'as grave tener Precision baja?
\end{classicbox}

\item ¿Qu\'e es el F1-Score y cu\'ando es m\'as \'util que el Accuracy?

\item Explique la curva ROC y el AUC-ROC. ¿Qu\'e significa un AUC de 0.5?

\item ¿Cu\'ando preferir\'ia la curva Precision-Recall sobre la curva ROC?

\end{enumerate}

% =============================================================
\section{Modelos Lineales}
% =============================================================

\subsection{Regresi\'on Lineal y Log\'istica}

\begin{enumerate}[label=\textbf{P\arabic*.}, resume]

\item ¿Por qu\'e la Regresi\'on Log\'istica es un modelo lineal si usa la funci\'on sigmoide?

\item ¿Por qu\'e la Regresi\'on Lineal tiene soluci\'on cerrada pero la Log\'istica no?

\end{enumerate}

\subsection{Regularizaci\'on}

\begin{enumerate}[label=\textbf{P\arabic*.}, resume]

\item \leavevmode

\begin{classicbox}
\textcolor{tdred}{\textbf{[CL\'ASICA]}} Compare L1 (Lasso) y L2 (Ridge). ¿En qu\'e
se diferencian y cu\'ando usar cada una?
\end{classicbox}

\item ¿Qu\'e efecto tiene aumentar $\lambda$ en la regularizaci\'on? ¿Y disminuirlo?

\end{enumerate}

% =============================================================
\section{\'Arboles de Decisi\'on}
% =============================================================

\begin{enumerate}[label=\textbf{P\arabic*.}, resume]

\item Explique intuitivamente por qu\'e los \'arboles de decisi\'on son modelos de
  \textbf{alta varianza}.

\item Defina Gini y Entrop\'ia. ¿Qu\'e valor toman cuando el nodo es puro? ¿Y cuando es
  m\'aximamente impuro?

\item \leavevmode

\begin{classicbox}
\textcolor{tdred}{\textbf{[CL\'ASICA]}} Mencione \textbf{tres} hiperpar\'ametros que
regularizan un \'arbol de decisi\'on (reducen overfitting) y explique el efecto de cada
uno.
\end{classicbox}

\item Un \'arbol fue entrenado con temperaturas entre 10\textdegree C y 40\textdegree C.
  En producci\'on recibe una observaci\'on de 55\textdegree C. ¿Qu\'e predice? ¿Es un
  problema?

\item ¿Qu\'e es la importancia de features en un \'arbol? ¿Cu\'al es su sesgo principal?

\end{enumerate}

% =============================================================
\section{Ensambles I: Bagging y Random Forest}
% =============================================================

\begin{enumerate}[label=\textbf{P\arabic*.}, resume]

\item Explique el concepto de \textbf{Bootstrap}. ¿Qu\'e porcentaje de muestras queda
  fuera (OOB) en promedio?

\item ¿C\'omo se usa el error OOB en Random Forest? ¿Qu\'e ventaja tiene frente a la
  validaci\'on cruzada?

\item \leavevmode

\begin{classicbox}
\textcolor{tdred}{\textbf{[CL\'ASICA]}} ¿Cu\'al es la diferencia fundamental entre
Bagging y Random Forest?
\end{classicbox}

\item ¿Agregar m\'as \'arboles a un Random Forest puede causar overfitting? ¿Y en
  Gradient Boosting?

\item ¿Por qu\'e la diversidad (decorrelaci\'on) entre modelos es importante en un
  ensamble?

\end{enumerate}

% =============================================================
\section{Ensambles II: Boosting}
% =============================================================

\begin{enumerate}[label=\textbf{P\arabic*.}, resume]

\item \leavevmode

\begin{classicbox}
\textcolor{tdred}{\textbf{[CL\'ASICA]}} Explique la diferencia conceptual entre Bagging
y Boosting en cuanto a c\'omo reducen el error.
\end{classicbox}

\item Explique el algoritmo de AdaBoost. ¿C\'omo se actualizan los pesos?

\item ¿Qu\'e son los pseudo-residuos en Gradient Boosting? ¿Por qu\'e se usan?

\item ¿Qu\'e rol cumplen el \textit{learning rate} ($\eta$) y el n\'umero de \'arboles
  ($M$) en Gradient Boosting? ¿C\'omo interact\'uan?

\item Compare XGBoost, LightGBM y CatBoost en al menos dos dimensiones.

\end{enumerate}

% =============================================================
\section{Reducci\'on de la Dimensionalidad}
% =============================================================

\begin{enumerate}[label=\textbf{P\arabic*.}, resume]

\item ¿Qu\'e es la maldici\'on de la dimensionalidad y por qu\'e afecta a los modelos
  basados en distancias?

\item Explique PCA paso a paso. ¿Qu\'e maximiza y c\'omo se obtienen las componentes?

\item ¿Qu\'e significa la varianza explicada en PCA? ¿C\'omo se elige $k$?

\item \leavevmode

\begin{classicbox}
\textcolor{tdred}{\textbf{[CL\'ASICA]}} Compare PCA, t-SNE y UMAP en cuanto a:
linealidad, velocidad, y si pueden proyectar nuevas muestras.
\end{classicbox}

\end{enumerate}

% =============================================================
\section{Overfitting, Underfitting y Trade-offs}
% =============================================================

\begin{enumerate}[label=\textbf{P\arabic*.}, resume]

\item \leavevmode

\begin{classicbox}
\textcolor{tdred}{\textbf{[CL\'ASICA --- SE TOMA SIEMPRE]}}
¿Qu\'e hiperpar\'ametros pueden causar \textbf{overfitting}? Liste al menos cinco
(de distintos modelos) y explique por qu\'e.
\end{classicbox}

\item Explique el Bias-Variance Trade-off con la f\'ormula y un ejemplo.

\item ¿C\'omo detectar si un modelo tiene overfitting o underfitting a partir de las
  curvas de aprendizaje?

\end{enumerate}

% =============================================================
\section{Selecci\'on de Features}
% =============================================================

\begin{enumerate}[label=\textbf{P\arabic*.}, resume]

\item Compare los tres tipos de selecci\'on de features: Filter, Wrapper y Embedded.

\item ¿Por qu\'e Label Encoding puede ser problem\'atico para modelos lineales?

\end{enumerate}

% =============================================================
\section{Resumen: Comparaci\'on de Modelos}
% =============================================================

\begin{enumerate}[label=\textbf{P\arabic*.}, resume]

\item \leavevmode

\begin{classicbox}
\textcolor{tdred}{\textbf{[CL\'ASICA]}} ¿Cu\'ando elegiría un \'arbol simple sobre
Random Forest o Gradient Boosting?
\end{classicbox}

\item Complete la siguiente tabla comparativa de modelos:

\vspace{0.3cm}
\begin{center}
\begin{tabular}{lcccc}
\toprule
\textbf{Modelo} & \textbf{Sesgo} & \textbf{Varianza} & \textbf{Interpretable} & \textbf{No lineal}\\
\midrule
Reg.\ Lineal/Log\'istica & & & &\\
\'Arbol (profundo) & & & &\\
\'Arbol (podado) & & & &\\
Bagging / RF & & & &\\
Gradient Boosting & & & &\\
\bottomrule
\end{tabular}
\end{center}

\end{enumerate}

% =============================================================
\section{Preguntas de Integraci\'on}
% =============================================================

\begin{enumerate}[label=\textbf{P\arabic*.}, resume]

\item Un equipo implementa el siguiente pipeline: (I) cargar dataset; (II) imputar con la
  media de todo el dataset; (III) escalar con MinMaxScaler sobre todo el dataset; (IV)
  dividir en train/test; (V) entrenar; (VI) evaluar. Identifique todos los errores.

\item Un modelo de detecci\'on de enfermedades tiene Recall = 0.99 pero Precision = 0.10.
  ¿C\'omo lo interpretar\'ia? ¿Es buen modelo?

\item Un banco quiere predecir impago de cr\'editos. El dataset tiene 500k clientes,
  solo 12k no pagaron. ¿Qu\'e problemas anticipar\'ia y c\'omo los abordar\'ia?

\end{enumerate}

\vspace{1cm}
\begin{center}
\rule{0.8\linewidth}{1pt}\\[6pt]
{\itshape Banco de preguntas completo --- TD VI 1S-2026}\\[4pt]
{\small\color{gray} Para las respuestas, consulte el documento \textit{preguntas\_teoricas.pdf}.}
\end{center}

\end{document}
